% Theorem A --- Q1-augmented Navier--Stokes (this PDE only).
% Not classical NS. Not Track B. Compile with pdflatex.
% Desk chain: docs/A-PROOF-CHAIN.md, docs/AUGMENTED-NS-PROOF-CHAIN.md
\documentclass[11pt]{article}
\usepackage[margin=1in]{geometry}
\usepackage{amsmath,amssymb,amsthm,mathtools}
\usepackage[T1]{fontenc}
\usepackage{microtype}
\usepackage{hyperref}

\newtheorem{theorem}{Theorem}
\newtheorem{lemma}{Lemma}
\newtheorem{remark}{Remark}

\title{Global regularity for the \(Q_1\)-augmented
Navier--Stokes system\\
\large Track A --- this PDE only}
\author{Jonathan Robert Simons\\
Prime Field Technologies}
\date{5 September 2026}

\begin{document}
\maketitle

\begin{abstract}
On the three-torus, the Navier--Stokes equations with an extra
Ladyzhenskaya / \(p\)-Laplacian stress
(\(Q_1\)-augmentation, \(\varepsilon>0\), \(\beta\ge 1/2\))
admit a unique global smooth solution. This is Theorem~A on
the Domain Architect desk. It is a theorem about
\textbf{this equation}. It is not classical three-dimensional
Navier--Stokes. The extra dissipation may be sent to zero
only after a separate uniform \(H^1\) bound, which remains
open; even then classical regularity is a further argument
(Track~B). The chain records a standard class
(Ladyzhenskaya 1968; M\'alek--Ne\v{c}as--R\r{u}\v{z}i\v{c}ka)
in \(Q_1\) notation.
\end{abstract}

\section{The system}

Let \(\nu>0\), \(\varepsilon>0\), \(\alpha>0\), and
\(\beta\ge 1/2\). On \(\mathbb{T}^3\times[0,\infty)\) consider
\begin{equation}
\label{eq:q1}
\partial_t u+(u\cdot\nabla)u
=-\nabla p+\nu\Delta u
+\varepsilon^\alpha\,\mathbb{P}\,
\operatorname{div}\bigl(|\nabla u|^\beta\nabla u\bigr),
\qquad
\nabla\cdot u=0,
\end{equation}
with \(u_0\in H^1(\mathbb{T}^3)\) divergence-free.
Here \(\mathbb{P}\) is the Leray projector.
Write \(p=\beta+2\), so \(\beta\ge 1/2\) is \(p\ge 5/2\).
The extra term is the variational derivative of
\(\frac1p\int|\nabla u|^p\). It is not the scalar
\(-\varepsilon^\alpha|\nabla u|^\beta\Delta u\).

\begin{theorem}[Theorem A]
\label{thm:A}
Let \(\nu>0\), \(\varepsilon>0\), \(\alpha>0\),
\(\beta\ge 1/2\), and let \(u_0\in H^1(\mathbb{T}^3)\) be
divergence-free. The system~\eqref{eq:q1} has a unique
solution
\[
u\in C^\infty(\mathbb{T}^3\times(0,\infty))
\cap L^\infty(0,\infty;H^1).
\]
No finite-time singularity occurs for this PDE.
The data need not be axisymmetric.
\end{theorem}

\section{Proof}

\begin{lemma}[Energy]
For a smooth solution of~\eqref{eq:q1},
\[
\frac12\frac{d}{dt}\|u\|_2^2
+\nu\|\nabla u\|_2^2
+\varepsilon^\alpha\|\nabla u\|_{L^{\beta+2}}^{\beta+2}
=0.
\]
Hence for every \(T>0\),
\[
\frac12\|u(T)\|_2^2
+\nu\int_0^T\|\nabla u\|_2^2\,dt
+\varepsilon^\alpha\int_0^T\|\nabla u\|_{L^{\beta+2}}^{\beta+2}\,dt
=\frac12\|u_0\|_2^2.
\]
\end{lemma}

\begin{proof}
Test~\eqref{eq:q1} against \(u\). The convective term and
the pressure vanish:
\(\int(u\cdot\nabla)u\cdot u=\frac12\int u\cdot\nabla(|u|^2)=0\).
The extra stress gives
\(\int\mathbb{P}\operatorname{div}(|\nabla u|^\beta\nabla u)\cdot u
=-\int|\nabla u|^{\beta+2}\).
The Stokes term is \(-\nu\|\nabla u\|_2^2\).
\end{proof}

\begin{lemma}[Galerkin]
The Galerkin projections onto the first \(n\) Stokes
eigenfunctions exist globally. The energy bounds are
uniform in \(n\). A subsequence converges weakly to a
weak solution
\[
u\in L^\infty(0,\infty;L^2)\cap L^2(0,\infty;H^1),
\qquad
\nabla u\in L^{\beta+2}(0,\infty;L^{\beta+2}).
\]
\end{lemma}

\begin{proof}
Finite-dimensional ODE; the same energy identity holds.
No finite-time blowup of \(\|u_n\|_2\). Banach--Alaoglu
and Aubin--Lions pass to the limit. Monotonicity of
\(v\mapsto\operatorname{div}(|\nabla v|^\beta\nabla v)\)
is Minty--Browder. Details:
M\'alek--Ne\v{c}as--R\r{u}\v{z}i\v{c}ka,
\textit{Weak and Measure-valued Solutions to Evolutionary PDEs},
Chapter~5.
\end{proof}

\begin{lemma}[Strong solutions, \(\beta\ge 1/2\)]
The weak solution is unique and lies in
\(L^\infty(0,\infty;H^1)\cap L^2(0,\infty;H^2)\).
The constant depends on
\(\varepsilon,\alpha,\beta,\nu,\|u_0\|_{H^1}\)
and blows up as \(\varepsilon\to 0\).
\end{lemma}

\begin{proof}
For \(\beta\ge 1/2\) one has \(\beta+2\ge 5/2\). Extra
integrability of \(\nabla u\) meets the Ladyzhenskaya
\(p\ge 5/2\) criterion (equivalently a
Ladyzhenskaya--Prodi--Serrin pair built from that
integrability). The extra term produces a monotone
remainder that absorbs the convective difference of two
solutions. See Ladyzhenskaya (1968) and
M\'alek--Ne\v{c}as--R\r{u}\v{z}i\v{c}ka.
This is extra integrability of \(\nabla u\), not a
geometric cancel, and not a \(\Phi\) identity.
\end{proof}

\begin{lemma}[Bootstrap]
The unique strong solution is
\(C^\infty(\mathbb{T}^3\times(0,\infty))\).
\end{lemma}

\begin{proof}
Frozen \(\varepsilon>0\), the linearized operator is a
uniformly elliptic Stokes operator. Difference quotients
give \(H^k\) for all \(k\). Sobolev embedding yields
spatial smoothness; time regularity follows from the
equation.
\end{proof}

\begin{proof}[Proof of Theorem~\ref{thm:A}]
Galerkin existence, weak limit, strong uniqueness at
\(\beta\ge 1/2\), then the bootstrap.
\end{proof}

\section{What this does not prove}

\begin{remark}
Uniform \(H^1\) as \(\varepsilon\to 0\) is not claimed.
A decaying \(Q_1\) integral is not that bound.
\end{remark}

\begin{remark}
Classical unaugmented Navier--Stokes (Track~B: keep
\(1/r^4\), no \(Q_1\)) is not claimed.
\(A\Rightarrow B\) remains false on this desk.
\end{remark}

\begin{remark}
Riemann zeros, Birch--Swinnerton-Dyer, Hodge, Yang--Mills,
and \(\mathrm{P}\) vs \(\mathrm{NP}\) are different objects.
\end{remark}

\begin{thebibliography}{9}
\bibitem{L68}
O.~A. Ladyzhenskaya,
\textit{The Mathematical Theory of Viscous Incompressible Flow},
Gordon and Breach, 1969
(modified NS: extra \(p\)-Laplacian stress).
\bibitem{MNR}
J.~M\'alek, J.~Ne\v{c}as, M.~R\r{u}\v{z}i\v{c}ka,
\textit{Weak and Measure-valued Solutions to Evolutionary PDEs},
Chapman \& Hall, 1996.
\end{thebibliography}

\end{document}
